
\documentclass{article} %210 mm × 297 mm
\usepackage[utf8]{inputenc}
\usepackage[a4paper,left=1.5cm, right=1.5cm, top=2cm, bottom=2cm]{geometry}
\usepackage{graphicx}
\usepackage{systeme}          % for Solution 3
\usepackage{array}  
%\usepackage{blindtext}
\usepackage{multirow}
\usepackage[normalem]{ulem}
\usepackage{mhchem}
\usepackage{url}
\usepackage{subfigure}
\usepackage{amsfonts,latexsym} % para tener disponibilidad de diversos simbolos
\usepackage{enumerate}
\usepackage{longdivision}
\usepackage{polynom}
%\usepackage{booktabs}
\usepackage{float}
%\usepackage{threeparttable}
\usepackage{array,colortbl}
%\usepackage{ifpdf}
%\usepackage{rotating}
\usepackage{stfloats}
\usepackage{url}
\usepackage{listings}
\usepackage{fancyhdr}
\usepackage{biblatex}
\usepackage{color}
\addbibresource{sources.bib}
%\usepackage{hyperref}
%\usepackage{wrapfig}
\usepackage{array}
%\usepackage{multirow}
%\usepackage{adjustbox}
%\usepackage{nccmath}
\usepackage[dvipsnames]{xcolor}
\usepackage{tcolorbox}
\newtcolorbox{mybox}{colback=white,
colframe=black}
\newcommand{\titulo}{Tutorium 9; Diskrete Mathematik}
\newcommand{\nombre}{Lil'Jana}
\newcommand{\FCE}{\textbf{WiSe 2024/2025}}

 

\usepackage{fancyhdr}
\pagestyle{fancy}
\fancyhead{}
\fancyfoot{}
\fancyhead[LO]{\small\nombre}
\fancyhead[RE]{\small\titulo}
\fancyhead[LO]{\small\FCE}
\fancyhead[RO]{\small 07.01.2025}
\fancyfoot[RO,LE] {\thepage}
\usepackage[onehalfspacing]{setspace}

\setcounter{page}{1} %Número de la página, la editorial lo cambiará
\begin{document}
\begin{tabular}{p{12cm}}
{{\Large\bf\titulo}}\\
\vspace{0.2cm}\\
 {\large\nombre}\\
 \vspace{0.1cm}
\end{tabular} \\
Inkl. ein paar als richtig befundene Ausschnitte aus meiner Abgabe. 
\section*{Besprochene Lösungen}

\subsection*{Aufgabe 1 (4 Punkte)}
Sei $S$ eine Differenzenfamilie in $\mathbb{Z}/m\mathbb{Z}$. Zeigen Sie, dass $S + i$ eine Differenzenfamilie in $\mathbb{Z}/m\mathbb{Z}$ ist für jedes $i \in \mathbb{Z}/m\mathbb{Z}$. Folgern Sie, dass man ohne Einschränkung annehmen kann, dass jede Differenzenfamilie $0$ und $1$ enthält.
\begin{mybox}
 Behauptung: Falls S eine Differenzenfamilie ist, dann ist S+1 eine Differenzenfamilie. \\
 Beweis: $T = S+i = \{x+i \mid x \in S\}$ \\
\begin{align*}
    \mid T_r \mid & = \mid \{(x,y) \in T \times T \mid x-y = r\}\mid  \\
    & = \mid \{ (u+i, v+i) \mid u,v \in S , \underbrace{(u+i) - (v+i) }_{u-v} = r \} \mid \\
    & = \mid \{ (u,v) \in S \times S \mid u-v = r \}\mid \\
     & = \mid S_r \mid 
\end{align*}
Somit ist $T_r$ unabhängig von $r \in \mathbb{Z}/m\mathbb{Z}$ und $T$ ist eine Differenzfamilie. \\ \newline
Falls $S_1 \neq leere Menge$, wähle $(\alpha, y) \in S_1$ \\
Dann $\{0,1\} \in S+ (-y)$
\end{mybox}

\subsection*{Aufgabe 2 (4 Punkte)}
Sei $S$ eine Differenzenfamilie in $\mathbb{Z}/m\mathbb{Z}$. Zeigen Sie, dass $T = \mathbb{Z}/m\mathbb{Z} \setminus S$ eine Differenzenfamilie ist. Bestimmen Sie die Parameter des von $T$ induzierten $2$-Designs in Abhängigkeit der Parameter des von $S$ induzierten $2$-Designs.
\begin{mybox}
Wie finden wir heraus, ob es eine Differenzenfamilie ist? Zählen. \\
Behauptung: Falls S eine Differenzenfamilie ist, dann ist $\mathbb{Z}/m\mathbb{Z} \setminus S$ eine Differenzenfamilie.  \\ \newline
\textbf{Beweis:} $T = \mathbb{Z}/m\mathbb{Z} \setminus S$ \\
$T_r = \{ (x,y)\}$ \\
$T_r = \{(x,y) \in T \times T \mid x-y = r\} $ \\
$ m = \mid \{(x,y) \in \mathbb{Z}/m\mathbb{Z} \times \mathbb{Z}/m\mathbb{Z} \mid x-y =r \}$ 
\begin{align*}
    \mid S_r\mid &= \mid \{ (x,y) \in S \times S \mid x-y = r \} \mid \\
    a_r: &= \mid \{ (x,y) \in S \times T \mid x-y = r \} \mid \\
    b_r: &= \mid \{ (x,y) \in T \times S \mid x-y = r \} \mid 
\end{align*}
\end{mybox}
\begin{mybox}
 \sout{Hauptfrage: Sind $a_r$ und $b_r$ unabhängig von $r$.  \\
Wie zeigen wir, dass zwei mengen die gleiche Anzahl von Elementen haben? Wir schreiben eine Bijektion hin. Also können wir versuchen eine Bijektion zwischen $a_r$ und $\mid \{ (x,y) \in S \times T \mid x-y = 1 \} \mid$ (die Menge wo wir r = 1 setzen) zu schreiben. }

Was können wir über $c_r = \mid \{ (x,y) \in S \times \mathbb{Z}/m\mathbb{Z}  \mid x-y = r \} \mid$ aussagen? \\
Dass $c_r := \mid \{ (x,y) \in S \times \mathbb{Z}/m\mathbb{Z}  \mid x-y = r \} \mid = \mid S \mid$. \\
ALso: $c_r = \mid S_r \mid + a_r $ \scriptsize{Weil weder die linke noch die mittlere Seite abhängig von r ist, muss $a_r$ unabhängig von r sein - nach dem gleichen Argument ist $b_r$ auch unabhängig von r (da muss man aber $\mathbb{Z}/m\mathbb{Z} \times S$ nehmen anstatt das oben) bzw. die Mächtigkeiten der beiden Mengen sind unabhängig von r} \normalsize \\
$d_r := \{(x,y) \in T \times \mathbb{Z}/m\mathbb{Z} \mid x-y =r\} \mid  = \mid T \mid = m - \mid S \mid$ \\
$m - \mid S \mid  = d_r = \mid T_r \mid + b_r$ \\
$\Rightarrow  \mid T_r \mid $ unabhängig von r \\ \newline
$S \subset \mathbb{Z}/m\mathbb{Z} $ Differnzenfamilie \\
$wobbelpfeil$ 2-Deisgn mit Parametern $v = m$, $k = \mid S\mid$, $r_2 = \frac{k(n-1)}{m-1}$ \\
$ T = \mathbb{Z}/m\mathbb{Z} \backslash S$ 
\begin{align*}
    v_T & = m = v \\
    k_T & = \mid T \mid  = m - \mid S \mid = m - k \\
    {(r_2)}_T & = \frac{k_T (k_T -1)}{m-1}
\end{align*}
\end{mybox}
\subsubsection*{Weiteres Beispiel:}
\begin{mybox}
 Bsp.: $S = \{1,2,3\} \subset \mathbb{Z}/7\mathbb{Z}$ \\
 $ T = \{0,3,5,6\}$\\
  Differenzentabelle: 
  \begin{table}[H]
      \centering
      \begin{tabular}{c|cccc}
      &0&3&5&6\\ \hline
           0&0&3&5&6  \\
           3&4&0&2&3 \\
           5&2&5&0&1 \\
           6&1&4&6&0
      \end{tabular}
  \end{table}
  Alles gleich oft da, ist also eine Differenzenfamilie. 
  
  \hrule
  
  Wenn r = 2 und x = 4 muss y = 2, weil 4-y = 2 \\
  Help wofür gilt das, das hat sie zu $T_r$ oben gesagt. 
\end{mybox}

\subsection*{Aufgabe 3 (4 Punkte)}
Ergänzen Sie die Menge von Codewörtern: $\{00000, \ 00110, \ 10101\}$ \\
zu einem linearen Code $C$. Finden Sie eine Kontrollmatrix für $C$ und bestimmen Sie das Gewicht aller Codewörter in $C$.
\begin{mybox}

    \begin{tabular}{c|cccc}
        & 00000 & 00110 & 10101 & 10011\\ \hline
   00000& 00000 & 00110 & 10101 & 10011\\
   00110&       & 00000 & 10011 & 10101\\
   10101&       &       & 00000 & 00110\\
   10011&       &       &       & 00000
    \end{tabular}
\end{mybox}
\begin{mybox}
Ergänzter Code: \{00000, \ 00110, \ 10101, \ 10011\}\\
Gewichte: 
\begin{align*}
    00000: & \quad 0\\
    00110: & \quad 2\\
    10101: & \quad 3\\
    10011: & \quad 3
\end{align*}
Basis: \{\ 00110, \ 10101\}, denn die beiden sind linear unabhängig (im Gegensatz zu 10011) und den Nullvektor bezieht man für eine Basis gar nicht mit ein. \\
Was zieht man zuerst in der Matrix? Zeilen oder Spalten? Zeilen. \\ 

Wir suchen eine Kontrollmatrix $H$ mit $_{ker}H = C$ \\
$dim\ C = 2$ \\
Wir können $H$ als $(3\times 5)$-Matrix wählen. 
\begin{equation*}
    A = \begin{pmatrix}
0 & 0 & 1 & 1 & 0 \\
1 & 0 & 1 & 0 & 1
\end{pmatrix}
\end{equation*}
\begin{equation*}
    _{ker}A = span \{ 
    \begin{pmatrix}
    1 & \\ 0 & \\ 1 & \\ 1 & \\ 0 
    \end{pmatrix}, 
    \begin{pmatrix}
    0 & \\ 1 & \\ 0 & \\ 0 & \\ 0
    \end{pmatrix}, 
    \begin{pmatrix}
    1 & \\ 0 & \\ 0 & \\ 0 & \\ 1
    \end{pmatrix}
    \}
\end{equation*}

\begin{equation*}
    H = \begin{pmatrix}
    1&0&1&1&0 \\
    0&1&0&0&0 \\
    1&0&0&0&1 
    \end{pmatrix}
\end{equation*} \\
\textbf{Anmerkungen}:\\
$_{ker}$ ist der Kern (Siehe Lineare Algebra)\\
$f: U \rightarrow V$ \\
$_{ker}f = \{u \mid f(u) = 0\}$ \\ \newline
$span$ ist der Spann, also die lineare Hülle 
\end{mybox}

\subsection*{Aufgabe 4 (4 Punkte)}
Der Code $C$ sei durch die folgende Kontrollmatrix gegeben:
\[
\begin{pmatrix}
1 & 0 & 1 & 0 & 1 & 1 \\
0 & 1 & 0 & 0 & 1 & 1 \\
1 & 0 & 0 & 1 & 1 & 0 \\
\end{pmatrix}
\]
\begin{enumerate}
    \item[(a)] Korrigieren Sie folgende empfangenen Wörter zu Codewörtern in $C$ unter der Annahme, dass höchstens ein Fehler aufgetreten ist:
    \[
    110111, \quad 101010, \quad 110101.
    \]
     \begin{mybox}
    \[
    110111: 
\begin{pmatrix} 
1 & 0 & 1 & 0 & 1 & 1 \\
0 & 1 & 0 & 0 & 1 & 1 \\
1 & 0 & 0 & 1 & 1 & 0
\end{pmatrix} 
\cdot 
\begin{pmatrix} 
1 \\
1 \\
0 \\
1 \\
1 \\
1
\end{pmatrix}
= \begin{pmatrix}
1 \cdot 1 + 0 \cdot 1 + 1 \cdot 0 + 0 \cdot 1 + 1 \cdot 1 + 1 \cdot 1 \\
0 \cdot 1 + 1 \cdot 1 + 0 \cdot 0 + 0 \cdot 1 + 1 \cdot 1 + 1 \cdot 1 \\
1 \cdot 1 + 0 \cdot 1 + 0 \cdot 0 + 1 \cdot 1 + 1 \cdot 1 + 0 \cdot 1
\end{pmatrix}
= \begin{pmatrix}
1 \\
1 \\
1
\end{pmatrix}
\]
$ 110111 \rightarrow 110101$
Wenn wir einen Fehler an der fünften Stelle von H haben, ist im empfangenen Wort das 5te Bit falsch!
\[
101010: 
\begin{pmatrix} 
1 & 0 & 1 & 0 & 1 & 1 \\
0 & 1 & 0 & 0 & 1 & 1 \\
1 & 0 & 0 & 1 & 1 & 0
\end{pmatrix}
\cdot 
\begin{pmatrix} 
1 \\
0 \\
1 \\
0 \\
1 \\
0
\end{pmatrix}
= 
\begin{pmatrix} 
1 \cdot 1 + 0 \cdot 0 + 1 \cdot 1 + 0 \cdot 0 + 1 \cdot 1 + 1 \cdot 0 \\
0 \cdot 1 + 1 \cdot 0 + 0 \cdot 1 + 0 \cdot 0 + 1 \cdot 1 + 1 \cdot 0 \\
1 \cdot 1 + 0 \cdot 0 + 0 \cdot 1 + 1 \cdot 0 + 1 \cdot 1 + 0 \cdot 0
\end{pmatrix}
= 
\begin{pmatrix} 
1 \\
1 \\
0
\end{pmatrix}
\]
$101010 \rightarrow 101011$
\[
110101: 
\begin{pmatrix} 
1 & 0 & 1 & 0 & 1 & 1 \\
0 & 1 & 0 & 0 & 1 & 1 \\
1 & 0 & 0 & 1 & 1 & 0
\end{pmatrix}
\cdot 
\begin{pmatrix} 
1 \\
1 \\
0 \\
1 \\
0 \\
1
\end{pmatrix}
= 
\begin{pmatrix} 
1 \cdot 1 + 0 \cdot 1 + 1 \cdot 0 + 0 \cdot 1 + 1 \cdot 0 + 1 \cdot 1 \\
0 \cdot 1 + 1 \cdot 1 + 0 \cdot 0 + 0 \cdot 1 + 1 \cdot 0 + 1 \cdot 1 \\
1 \cdot 1 + 0 \cdot 1 + 0 \cdot 0 + 1 \cdot 1 + 1 \cdot 0 + 0 \cdot 1
\end{pmatrix}
= 
\begin{pmatrix} 
0 \\
0 \\
0
\end{pmatrix}
\]
Bleibt so, ist richtig. 
    \end{mybox} 
    \item[(b)] Wie viele Wörter existieren, die nicht durch Änderung von höchstens einem Bit zu einem Codewort in $C$ korrigiert werden können?
\end{enumerate}
\begin{mybox}
 Anzahl Wörter: $2^6$ \\
 Dim vom Kern (dim C): 3 \\
 Anzahl Codewörter: $2^3$ \\
 rang von H: 3, weil wir in H drei Zeilen haben (die linear unabhängig sind...) \\
 dim C + rang H = Länge eines Wortes \\
 Hier: $3+3 =6$ - passt \\ \newline
 Zu jedem Codewort gibt es 6 Wörter, die zu dem Codewort korrigiert werden können. \\ \newline
Wörter, die nicht durch Änderung eines Bits zu einem Cidewort korrigiert werden können: $2^6 - 2^3 - 6\cdot 2^3 = 64 - 8 - 48 = 8$
\end{mybox} 

\subsection*{Polynomdivision Crashkurs}
\begin{mybox}
 \[
\frac{ x^5 + x^4 + 2x^3 + x^2 + 4x + 2 }{x^2 + 2x + 3} 
 \]
 Ihre schreibweise: $(x^5 + x^4 + 2x^3 + x^2 + 4x + 2) = (x^2 + 2x + 3) ( ...)$
 Man schaut sich zuerst den höchsten an: $x^5$ - womit muss man $x^3$ multiplizieren, damit man das bekommt ? $x^2$ \\
Usw. ... \\ \newline
Ihre schreibweise: $(x^5 + x^4 + 2x^3 + x^2 + 4x + 2) = (x^2 + 2x + 3) (x^3 + 4x^2 + x +2 ) + (2x +1)$
Hier einmal als normale Polynomdivision, wo am Ende modulo 5 gerechnet wird: \\
  Hier drunter ihre Tabelle/Bruchrechnung aus Notizheft... \\ \newline
\polylongdiv[style=A]{x^5 + x^4 + 2x^3 + x^2 + 4x + 2 }{x^2 + 2x + 3}
\[- 3x - 4 \mod 5 = 2x +1\]
Und hier einmal (weniger gut formatiert) die Schreibweise von der Tafel, wo direkt immer module 5 gerechnet wird:  \\ \newline 
$\begin{array}{cccccccc}
    & x^5 & +  2x^3 & + x^2 &+ 4x &+ 2\\
    -(&  x^5 & + 2x^4 & + 3x^3& )  \\ \hline
    & & 4x^4 & 4x^3 & + x^2 \\
    &  -( & 4x^4 & + 3x^3 & + 2x^2& ) \\ \hline
    & & & x^3 & + 4x^2 & + 4x \\
    & & -( & x^3 & + 2x^2 & + 3x& ) \\ \hline
    & & & & 2x^2 & + x & +1 \\
    & & & -( & 2x^2 & + 4x & +1 & ) \\ \hline
    & & & & & 2x & +1 
\end{array}$

\end{mybox}
\end{document}